\subsection*{The Picture}

\begin{remark*}[Exposition]
The ordinary construction of the integers in this volume uses formal
differences: an integer is represented by a pair of whole numbers, and equal
formal differences are identified by an equivalence relation. That construction
is canonical for the book because it explains exactly how subtraction is
adjoined to the whole numbers.
\end{remark*}

\begin{remark*}[Exposition]
There is another natural picture. Integers look like two one-way rays separated
by zero:
\[
\cdots,\;-3,\;-2,\;-1,\;0,\;1,\;2,\;3,\;\cdots .
\]
The question is whether the integers can be constructed directly as a positive
ray, a negative ray, and a distinguished zero, with successor and predecessor
moving along those rays. The answer is yes, but the construction makes
different proof obligations visible.
\end{remark*}

\begin{definition}[Two-sided integer rays]
\label{def:two-sided-integer-rays}
A \emph{two-sided ray presentation} consists of three disjoint parts:
\[
Z=\{0\}\sqcup P\sqcup N,
\]
where \(P\) is generated from a first positive element by repeated successor
and \(N\) is generated from a first negative element by repeated predecessor.
\end{definition}
\SourceVariantOf{thm:canonical-form-of-integers}{LRA}{Volume II, Constructing a Number System}{source_variant_of}

\begin{remark*}[Standard quantified statement]
The carrier is partitioned into zero, a positive ray, and a negative ray.
\end{remark*}

\begin{remark*}[Interpretation]
This starts from the sign picture that is usually drawn after the integers are
already known. In the present section it is treated as construction data.
\end{remark*}

\begin{dependencies}
\begin{itemize}
  \item \hyperref[def:number-system-construction-data]{Construction Data}
  \item \hyperref[thm:canonical-form-of-integers]{Canonical Form of Integers}
\end{itemize}
\end{dependencies}

\subsection*{Construction Data}

\begin{definition}[Two-sided successor construction data]
\label{def:two-sided-successor-construction-data}
The \emph{two-sided successor construction data} consists of a carrier \(Z\),
a distinguished element \(0\in Z\), and two total functions
\[
\mathsf{succ},\mathsf{pred}:Z\to Z
\]
such that
\[
\mathsf{pred}(\mathsf{succ}(x))=x
\quad\text{and}\quad
\mathsf{succ}(\mathsf{pred}(x))=x
\]
for every \(x\in Z\), together with the decomposition
\[
Z=\{0\}\sqcup P\sqcup N
\]
where \(P\) is generated from \(\mathsf{succ}(0)\) by successor and
\(N\) is generated from \(\mathsf{pred}(0)\) by predecessor.
\end{definition}
\SourceVariantOf{def:integers}{Iwanus-Wybraniec}{Polish integer arithmetic axiomatizations, 1965}{source_variant_of}

\begin{remark*}[Standard quantified statement]
The raw data are zero, mutually inverse successor and predecessor maps, and a
disjoint exhaustion by the zero point, the positive successor ray, and the
negative predecessor ray.
\end{remark*}

\begin{remark*}[Interpretation]
The quotient construction makes equality the first problem. The two-sided
construction makes case analysis the first problem. Every operation must say
what happens at zero, on the positive ray, and on the negative ray.
\end{remark*}

\begin{remark*}[Exposition]
Lean verification record. Source: Iwanu{\'s}, \emph{Some Axiomatic Systems of
Integers Arithmetic}; Wybraniec, \emph{A Certain System of Axioms in the
Integer Arithmetic}. Formalizes
\hyperref[def:two-sided-successor-construction-data]{two-sided successor
construction data}. Lean module:
\texttt{LRA.VolumeII.Integers.Polish.TwoSidedSuccessor}. Declarations:
\texttt{Z}, \texttt{succ}, \texttt{pred}. Status: checked against
\texttt{lra-lean}.
\end{remark*}
\LeanFormalizes{def:two-sided-successor-construction-data}{lra-lean}{LRA.VolumeII.Integers.Polish.TwoSidedSuccessor}{Z}{checked}
\LeanFormalizes{def:two-sided-successor-construction-data}{lra-lean}{LRA.VolumeII.Integers.Polish.TwoSidedSuccessor}{succ}{checked}
\LeanFormalizes{def:two-sided-successor-construction-data}{lra-lean}{LRA.VolumeII.Integers.Polish.TwoSidedSuccessor}{pred}{checked}

\begin{dependencies}
\begin{itemize}
  \item \hyperref[def:number-system-construction-data]{Construction Data}
  \item \hyperref[def:integers]{Integers}
\end{itemize}
\end{dependencies}

\subsection*{Axioms and Induction}

\begin{theorem}[Successor and predecessor are inverse operations]
\label{thm:two-sided-successor-predecessor-inverses}
In the two-sided construction, successor and predecessor satisfy
\[
\mathsf{pred}(\mathsf{succ}(x))=x
\quad\text{and}\quad
\mathsf{succ}(\mathsf{pred}(x))=x
\]
for every \(x\in Z\). Consequently both maps are injective.
\hyperref[prf:two-sided-successor-predecessor-inverses]{Go to proof.}
\end{theorem}

\begin{remark*}[Standard quantified statement]
Successor and predecessor undo each other everywhere on \(Z\).
\end{remark*}

\begin{remark*}[Interpretation]
This replaces the one-sided Peano successor with a two-sided motion. Moving one
step right and then one step left returns to the same integer, and conversely.
\end{remark*}

\begin{remark*}[Exposition]
Lean verification record. Source: Iwanu{\'s}; Wybraniec. Formalizes
\hyperref[thm:two-sided-successor-predecessor-inverses]{the inverse-operation
theorem}. Lean module:
\texttt{LRA.VolumeII.Integers.Polish.TwoSidedSuccessor}. Declarations:
\texttt{pred\_succ}, \texttt{succ\_pred}, \texttt{succ\_injective},
\texttt{pred\_injective}. Status: checked against \texttt{lra-lean}.
\end{remark*}
\LeanFormalizes{thm:two-sided-successor-predecessor-inverses}{lra-lean}{LRA.VolumeII.Integers.Polish.TwoSidedSuccessor}{pred_succ}{checked}
\LeanFormalizes{thm:two-sided-successor-predecessor-inverses}{lra-lean}{LRA.VolumeII.Integers.Polish.TwoSidedSuccessor}{succ_pred}{checked}
\LeanFormalizes{thm:two-sided-successor-predecessor-inverses}{lra-lean}{LRA.VolumeII.Integers.Polish.TwoSidedSuccessor}{succ_injective}{checked}
\LeanFormalizes{thm:two-sided-successor-predecessor-inverses}{lra-lean}{LRA.VolumeII.Integers.Polish.TwoSidedSuccessor}{pred_injective}{checked}

\begin{dependencies}
\begin{itemize}
  \item \hyperref[def:two-sided-successor-construction-data]{Two-sided successor construction data}
\end{itemize}
\end{dependencies}

\begin{theorem}[Two-sided induction and recursion]
\label{thm:two-sided-induction-and-recursion}
Let \(A\) be a predicate on \(Z\). If \(A(0)\), \(A(x)\Rightarrow
A(\mathsf{succ}(x))\), and \(A(x)\Rightarrow A(\mathsf{pred}(x))\)
for every \(x\in Z\), then \(A(x)\) holds for every \(x\in Z\).
Moreover, any value at zero and mutually inverse successor and predecessor
steps determine a unique recursive function out of \(Z\).
\hyperref[prf:two-sided-induction-and-recursion]{Go to proof.}
\end{theorem}
\SourceVariantOf{def:number-system-construction-verification}{Iwanus-Wybraniec}{Polish integer arithmetic axiomatizations, 1965}{reduces_to}

\begin{remark*}[Standard quantified statement]
A property holds everywhere on \(Z\) once it holds at zero and propagates both
forward and backward; a recursive definition on \(Z\) is determined by its
zero value and compatible forward and backward steps.
\end{remark*}

\begin{remark*}[Interpretation]
This is why the ray decomposition matters. To prove something for every
integer, prove it at zero, along the positive ray, and along the negative ray.
To define an operation recursively, the forward and backward clauses must
agree at the places where the rays meet zero.
\end{remark*}

\begin{remark*}[Exposition]
Lean verification record. Source: Iwanu{\'s}; Wybraniec. Formalizes
\hyperref[thm:two-sided-induction-and-recursion]{two-sided induction and
recursion}. Lean module:
\texttt{LRA.VolumeII.Integers.Polish.TwoSidedSuccessor}. Declarations:
\texttt{twoSidedInduction}, \texttt{recZ}, \texttt{recursion\_exists},
\texttt{recursion\_unique}. Status: checked against \texttt{lra-lean}.
\end{remark*}
\LeanFormalizes{thm:two-sided-induction-and-recursion}{lra-lean}{LRA.VolumeII.Integers.Polish.TwoSidedSuccessor}{twoSidedInduction}{checked}
\LeanFormalizes{thm:two-sided-induction-and-recursion}{lra-lean}{LRA.VolumeII.Integers.Polish.TwoSidedSuccessor}{recZ}{checked}
\LeanFormalizes{thm:two-sided-induction-and-recursion}{lra-lean}{LRA.VolumeII.Integers.Polish.TwoSidedSuccessor}{recursion_exists}{checked}
\LeanFormalizes{thm:two-sided-induction-and-recursion}{lra-lean}{LRA.VolumeII.Integers.Polish.TwoSidedSuccessor}{recursion_unique}{checked}

\begin{dependencies}
\begin{itemize}
  \item \hyperref[thm:two-sided-successor-predecessor-inverses]{Successor and predecessor are inverse operations}
  \item \hyperref[def:number-system-construction-verification]{Construction Verification}
\end{itemize}
\end{dependencies}

\subsection*{Arithmetic from Motion}

\begin{remark*}[Exposition]
Addition is defined by moving the left input along the ray described by the
right input: adding zero does nothing, adding a positive element applies
successor repeatedly, and adding a negative element applies predecessor
repeatedly. This is conceptually simple, but it creates a long lemma ladder.
One must prove that successor commutes with addition, predecessor commutes with
addition, zero is a left identity, addition is commutative, and addition is
associative.
\end{remark*}

\begin{remark*}[Why the lemmas are needed]
The quotient construction pays for arithmetic through well-definedness. The
two-sided construction pays through structural rewriting. For example,
commutativity of addition is not built into the definition: the definition
recurses on the second input, so proving \(x+y=y+x\) requires separate
successor and predecessor transport lemmas.
\end{remark*}

\begin{theorem}[The two-sided construction gives an ordered commutative ring]
\label{thm:two-sided-integers-ordered-commutative-ring}
With addition, negation, multiplication, and order defined by the two-sided
successor construction, the resulting system satisfies the laws of an ordered
commutative ring.
\hyperref[prf:two-sided-integers-ordered-commutative-ring]{Go to proof.}
\end{theorem}
\SourceVariantOf{thm:z-commutative-ring}{Iwanus-Wybraniec}{Polish integer arithmetic axiomatizations, 1965}{reduces_to}
\SourceVariantOf{thm:z-totally-ordered-set}{Iwanus-Wybraniec}{Polish integer arithmetic axiomatizations, 1965}{reduces_to}

\begin{remark*}[Standard quantified statement]
The two-sided system has associative and commutative addition, additive
inverses, associative and commutative multiplication, distributivity, a total
order, and the usual compatibility of order with addition and multiplication
by positive elements.
\end{remark*}

\begin{remark*}[Interpretation]
The construction is not shorter than the formal-difference construction. Its
pedagogical value is that every law can be traced to motion along the two rays.
\end{remark*}

\begin{remark*}[Exposition]
Lean verification record. Source: Iwanu{\'s}; Wybraniec. Formalizes
\hyperref[thm:two-sided-integers-ordered-commutative-ring]{the ordered-ring
workup}. Lean modules:
\texttt{LRA.VolumeII.Integers.Polish.LandauWorkup} and
\texttt{LRA.VolumeII.Integers.Polish.Implementation}. Declarations:
\texttt{add\_comm}, \texttt{add\_assoc}, \texttt{mul\_comm},
\texttt{mul\_assoc}, \texttt{distrib\_left}, \texttt{distrib\_right},
\texttt{lt\_trichotomy}, \texttt{PolishOrderedRingLaws}. Status: checked
against \texttt{lra-lean}.
\end{remark*}
\LeanFormalizes{thm:two-sided-integers-ordered-commutative-ring}{lra-lean}{LRA.VolumeII.Integers.Polish.LandauWorkup}{add_comm}{checked}
\LeanFormalizes{thm:two-sided-integers-ordered-commutative-ring}{lra-lean}{LRA.VolumeII.Integers.Polish.LandauWorkup}{add_assoc}{checked}
\LeanFormalizes{thm:two-sided-integers-ordered-commutative-ring}{lra-lean}{LRA.VolumeII.Integers.Polish.LandauWorkup}{mul_comm}{checked}
\LeanFormalizes{thm:two-sided-integers-ordered-commutative-ring}{lra-lean}{LRA.VolumeII.Integers.Polish.LandauWorkup}{mul_assoc}{checked}
\LeanFormalizes{thm:two-sided-integers-ordered-commutative-ring}{lra-lean}{LRA.VolumeII.Integers.Polish.LandauWorkup}{distrib_left}{checked}
\LeanFormalizes{thm:two-sided-integers-ordered-commutative-ring}{lra-lean}{LRA.VolumeII.Integers.Polish.LandauWorkup}{distrib_right}{checked}
\LeanFormalizes{thm:two-sided-integers-ordered-commutative-ring}{lra-lean}{LRA.VolumeII.Integers.Polish.LandauWorkup}{lt_trichotomy}{checked}
\LeanFormalizes{thm:two-sided-integers-ordered-commutative-ring}{lra-lean}{LRA.VolumeII.Integers.Polish.Implementation}{PolishOrderedRingLaws}{checked}

\begin{dependencies}
\begin{itemize}
  \item \hyperref[thm:two-sided-induction-and-recursion]{Two-sided induction and recursion}
  \item \hyperref[thm:z-commutative-ring]{\(\mathbb{Z}\) Is a Commutative Ring}
  \item \hyperref[thm:z-totally-ordered-set]{\(\mathbb{Z}\) Is a Totally Ordered Set}
\end{itemize}
\end{dependencies}

\subsection*{Comparison with the Standard Construction}

\begin{theorem}[The two-sided integers are isomorphic to the usual integers]
\label{thm:two-sided-integers-isomorphic-to-z}
The two-sided integer construction is isomorphic to the usual integer system.
The isomorphism preserves zero, successor, predecessor, addition,
multiplication, negation, and order.
\hyperref[prf:two-sided-integers-isomorphic-to-z]{Go to proof.}
\end{theorem}
\SourceVariantOf{def:integers}{Iwanus-Wybraniec}{Polish integer arithmetic axiomatizations, 1965}{reduces_to}

\begin{remark*}[Standard quantified statement]
There is a structure-preserving bijection between the two-sided construction
and the canonical integer system.
\end{remark*}

\begin{remark*}[Interpretation]
This is the comparison theorem that prevents the construction from becoming a
new arithmetic universe. It is a different presentation of the same integers.
\end{remark*}

\begin{remark*}[Exposition]
Lean verification record. Source: Iwanu{\'s}; Wybraniec. Formalizes
\hyperref[thm:two-sided-integers-isomorphic-to-z]{the comparison theorem}.
Lean modules: \texttt{LRA.VolumeII.Integers.Polish.LandauWorkup} and
\texttt{LRA.VolumeII.Integers.Polish.DiscreteIntegerStructure}. Declarations:
\texttt{Z\_iso\_Int}, \texttt{categoricity}. Status: checked against
\texttt{lra-lean}.
\end{remark*}
\LeanFormalizes{thm:two-sided-integers-isomorphic-to-z}{lra-lean}{LRA.VolumeII.Integers.Polish.LandauWorkup}{Z_iso_Int}{checked}
\LeanFormalizes{thm:two-sided-integers-isomorphic-to-z}{lra-lean}{LRA.VolumeII.Integers.Polish.DiscreteIntegerStructure}{categoricity}{checked}

\begin{remark*}[Exposition]
The formal-difference construction begins with pairs and an equivalence
relation. Tao emphasizes expressions \(a{--}b\); Mendelson emphasizes
equivalence classes of pairs and develops more abstract algebra around them.
The two-sided construction begins instead with the geometry of the integer
line. It avoids quotient well-definedness, but it requires explicit recursive
definitions and many transport lemmas for successor and predecessor.
\end{remark*}

\begin{remark*}[Exposition]
This section is therefore not a replacement for the standard construction. It
is a worked example of construction verification: choose a target, propose
data, state the axioms, build operations, prove the laws, and finally compare
the result with the already accepted system.
\end{remark*}

\begin{dependencies}
\begin{itemize}
  \item \hyperref[def:two-sided-successor-construction-data]{Two-sided successor construction data}
  \item \hyperref[thm:two-sided-integers-ordered-commutative-ring]{The two-sided construction gives an ordered commutative ring}
  \item \hyperref[def:integers]{Integers}
\end{itemize}
\end{dependencies}
